\documentclass[../DoAn.tex]{subfiles}
\begin{document}

Chương 6 tổng kết những kết quả chính đạt được trong quá trình thực hiện đề tài tốt nghiệp ``Xây dựng phần mềm tối ưu hóa quy trình điểm danh và quản lý học sinh''. Trên cơ sở đó, chương cũng chỉ ra các hạn chế còn tồn tại và đề xuất những hướng phát triển tiếp theo để hệ thống có thể hoàn thiện hơn khi triển khai trong môi trường giáo dục thực tế.

\section{Kết luận}
\label{section:6.1}

\subsection{So sánh kết quả nghiên cứu với các giải pháp tương tự}
Kết quả của đồ án được đối chiếu với một số phương pháp điểm danh và giám sát học tập phổ biến nhằm làm rõ giá trị ứng dụng của hệ thống:
\begin{itemize}
    \item \textbf{So với điểm danh truyền thống}: Cách gọi tên hoặc ký tên giấy thường mất nhiều thời gian, đặc biệt ở các lớp đông từ $50$ đến $100$ sinh viên, đồng thời dễ phát sinh sai sót hoặc điểm danh hộ. Giải pháp Group FaceID của đồ án cho phép giảng viên chụp ảnh tập thể lớp học, sau đó hệ thống tự động phát hiện khuôn mặt, trích xuất đặc trưng ArcFace \cite{Deng_2019_CVPR} và ghi nhận kết quả trong khoảng $1.5$ đến $2.5$ giây.
    \item \textbf{So với nhận diện sinh trắc học cá nhân}: Các mô hình quét vân tay hoặc FaceID từng người yêu cầu sinh viên xếp hàng lần lượt, dễ gây ùn tắc trước giờ học. Hệ thống của đồ án xử lý đồng thời nhiều khuôn mặt trong một ảnh nhờ Greedy Matching và Lazy Cropping, từ đó giảm thao tác cho cả giảng viên lẫn sinh viên.
    \item \textbf{So với các phần mềm AI Proctoring thương mại}: Nhiều giải pháp giám sát thi trực tuyến cần truyền video liên tục lên máy chủ và xử lý bằng mô hình nặng, làm tăng chi phí GPU và băng thông. Đồ án lựa chọn hướng kết hợp điểm mốc khuôn mặt 2D trên trình duyệt với solvePnP để ước lượng tư thế đầu, giúp hệ thống có thể vận hành trên CPU thông thường với độ trễ WebSocket dưới $80$ ms.
    \item \textbf{So với quản lý dữ liệu bằng bảng tính thủ công}: Cách tổng hợp chuyên cần, điểm số và sự kiện thi bằng Excel dễ phân tán dữ liệu, khó đối soát khi có khiếu nại. Hệ thống tập trung hóa dữ liệu theo lớp học phần, cung cấp dashboard, lịch sử chuyên cần, bảng điểm và nhật ký giám sát để giảng viên, học viên và quản trị viên theo dõi minh bạch hơn.
\end{itemize}

\subsection{Các kết quả cụ thể đạt được của đề tài}
Trong phạm vi đồ án, hệ thống đã hoàn thành các mục tiêu chính đã đặt ra:
\begin{itemize}
    \item \textbf{Kiến trúc hệ thống}: Xây dựng mô hình Microservices gồm backend Spring Boot, dịch vụ AI nhận diện khuôn mặt nhóm bằng FastAPI, dịch vụ giám sát thi trực tuyến bằng FastAPI \& WebSockets và hệ thống xác thực tập trung Keycloak. Toàn bộ các thành phần được container hóa và triển khai qua Docker Compose trên VPS Ubuntu.
    \item \textbf{Ứng dụng người dùng}: Hoàn thiện ứng dụng di động React Native cho học sinh với các chức năng check-in FaceID, GPS/OTP, xem lịch sử điểm danh và làm bài kiểm tra. Đồng thời, xây dựng ứng dụng web React cho giảng viên/quản trị viên để tạo ca học, quản lý lớp, xem báo cáo và theo dõi thi trực tuyến.
    \item \textbf{Tích hợp AI và thuật toán}: Tích hợp các thành phần RetinaFace \cite{Deng_2020_CVPR}, ArcFace \cite{Deng_2019_CVPR}, MediaPipe FaceMesh \cite{Lugaresi2019MediaPipe}, solvePnP và Haversine vào các nghiệp vụ nhận diện, giám sát và định vị. Các kỹ thuật này được kết hợp theo hướng phục vụ trực tiếp cho quy trình điểm danh thay vì chỉ dừng ở mức thử nghiệm mô hình.
    \item \textbf{Kiểm thử thực tế}: Hệ thống được thử nghiệm trong môi trường giảng đường với $150$ học viên và $10$ giảng viên. Kết quả ghi nhận hệ thống hoạt động ổn định, độ trễ API dưới $150$ ms, cơ chế đối soát chuyên cần hoạt động hiệu quả qua việc kết hợp GPS/FaceID và dữ liệu chuyên cần, điểm số, nhật ký giám sát được hiển thị thống nhất trên các giao diện tra cứu.
\end{itemize}

\subsection{Các mặt hạn chế còn tồn tại}
Bên cạnh các kết quả đã đạt được, hệ thống vẫn còn một số hạn chế cần tiếp tục cải thiện:
\begin{itemize}
    \item \textbf{Nhận diện ảnh nhóm}: Với lớp quá đông hoặc phòng học sâu, khuôn mặt ở hàng cuối có thể nhỏ, mờ hoặc bị che khuất, làm giảm chất lượng phát hiện của RetinaFace và độ tin cậy khi so khớp bằng ArcFace.
    \item \textbf{Độ ổn định của điểm danh theo vị trí}: Việc xác nhận vị trí phụ thuộc vào chất lượng tín hiệu định vị của điện thoại. Trong phòng kín, tầng hầm hoặc khu vực nhiều vật cản, sai số GPS có thể tăng cao và gây từ chối check-in nhầm.
    \item \textbf{Dữ liệu khuôn mẫu}: Ảnh mẫu hiện được lưu tại thời điểm đăng ký, nên hệ thống có thể giảm độ chính xác khi học sinh thay đổi diện mạo đáng kể nếu không cập nhật lại dữ liệu.
    \item \textbf{Bằng chứng giám sát thi}: Luồng WebSocket phụ thuộc vào kết nối mạng của thí sinh. Khi mạng bị ngắt quãng, dữ liệu trong \path{gaze_buffer} có thể bị thiếu khung hình, làm video bằng chứng chưa thật sự liền mạch.
\end{itemize}

\subsection{Các bài học kinh nghiệm rút ra}
Quá trình thực hiện đồ án giúp tác giả rút ra một số kinh nghiệm quan trọng:
\begin{itemize}
    \item \textbf{Tư duy thiết kế hệ thống}: Một giải pháp AI hiệu quả không chỉ phụ thuộc vào mô hình, mà còn nằm ở cách tích hợp vào luồng nghiệp vụ. Các kỹ thuật như padding hộp giới hạn, Lazy Cropping và solvePnP giúp cân bằng giữa độ chính xác, tốc độ xử lý và chi phí hạ tầng.
    \item \textbf{Kỹ năng triển khai thực tế}: Đề tài giúp tác giả hiểu rõ hơn về Microservices, WebSockets, xử lý bất đồng bộ, bảo mật API qua Keycloak SSO và các cơ chế kiểm tra phần cứng di động để bảo vệ tính toàn vẹn dữ liệu điểm danh.
\end{itemize}

\section{Hướng phát triển}
\label{section:6.2}

Để hệ thống có thể vận hành ổn định hơn và mở rộng cho nhiều cơ sở đào tạo, các hướng phát triển tiếp theo được chia thành nhóm ngắn hạn và dài hạn.

\subsection{Các công việc cần thiết hoàn thiện hệ thống trong ngắn hạn}
Trong giai đoạn gần, hệ thống cần tập trung vào các cải tiến có thể nâng cao trực tiếp trải nghiệm sử dụng:
\begin{itemize}
    \item \textbf{Cải tiến điểm danh nhóm}: Bổ sung cơ chế chụp nhiều ảnh ở các góc khác nhau hoặc áp dụng Image Stitching để tạo ảnh lớp học rõ hơn, từ đó giảm lỗi với học sinh ngồi xa hoặc bị che khuất.
    \item \textbf{Tăng độ tin cậy của điểm danh trong nhà}: Kết hợp thêm BSSID của Wi-Fi phòng học hoặc tín hiệu BLE Beacon để hỗ trợ check-in trong nhà, nhất là ở những khu vực GPS yếu.
    \item \textbf{Nâng cấp giao diện giám sát thi}: Bổ sung cảnh báo âm thanh, popup thời gian thực và khả năng xem nhanh luồng webcam của thí sinh bị nghi ngờ gian lận, giúp giám thị phản ứng kịp thời hơn.
\end{itemize}

\subsection{Các định hướng phát triển mở rộng trong dài hạn}
Về lâu dài, hệ thống có thể được phát triển theo các hướng sau:
\begin{itemize}
    \item \textbf{Học gia tăng cho dữ liệu khuôn mặt}: Cho phép hệ thống cập nhật dần vector đặc trưng khi học sinh điểm danh thành công với độ tin cậy cao, giúp dữ liệu khuôn mẫu thích ứng tốt hơn với sự thay đổi diện mạo tự nhiên.
    \item \textbf{Giám sát thi đa chế độ}: Mở rộng phân hệ \path{ai-proctor} bằng giám sát âm thanh, phát hiện chuyển tab trình duyệt và kiểm tra thiết bị ngoại vi, từ đó tăng khả năng phát hiện gian lận trong thi trực tuyến.
    \item \textbf{Mô hình SaaS đa đơn vị}: Tái cấu trúc cơ sở dữ liệu và phân hệ quản trị theo hướng Multi-tenant, cho phép nhiều trường học hoặc trung tâm giáo dục sử dụng chung hạ tầng nhưng vẫn tách biệt dữ liệu.
\end{itemize}

Tổng kết lại, đồ án đã xây dựng được một hệ thống có khả năng hỗ trợ điểm danh và quản lý học sinh theo hướng tự động, nhanh và có tính bảo mật cao hơn so với phương pháp thủ công. Dù vẫn còn các hạn chế về ảnh nhóm, tín hiệu định vị và độ ổn định của luồng giám sát thi, các kết quả đạt được cho thấy hệ thống có tiềm năng tiếp tục hoàn thiện để triển khai rộng rãi trong môi trường giáo dục.

\end{document}
